\section{Cài đặt, giới hạn và hướng phát triển}
\subsection{Cách chạy}
Sau khi cài đặt theo README, nhóm em khởi động chương trình bằng:
\begin{verbatim}
npm install
npm run dev
\end{verbatim}
Trong giao diện, người dùng chọn dataset, thuật toán, scenario, start và goal, sau đó bấm \textit{Run Algorithm}. Với bài toán nhiều điểm, chọn depot và thêm các stop trước khi chạy phần tối ưu tour. Các lệnh kiểm thử và benchmark được giữ trong README để nhóm có thể kiểm tra lại artifact khi cần.

\subsection{Giới hạn của đồ án}
Dataset hiện tại chưa phải dữ liệu traffic real-time. Một phần traffic/flood trong fixture là dữ liệu \texttt{SIMULATED}; một số nhận định cho dữ liệu tương lai vẫn là \texttt{ASSUMPTION}. Nếu cho phép người dùng nhập custom weights thì cần kiểm tra giá trị đầu vào và thang đo. Held-Karp có độ phức tạp tăng theo cấp số mũ nên chỉ phù hợp với số stop nhỏ. Heuristic lower-bound bảo đảm tính tối ưu cho composite cost hiện tại, nhưng có thể không mạnh nếu cost profile thay đổi nhiều.

\subsection{Hướng phát triển}
Trong tương lai, nhóm em có thể bổ sung traffic real-time, thêm scenario peak/rain cho landmark release, hỗ trợ custom weights có validation, bài toán nhiều xe và time windows. Giao diện cũng có thể được mở rộng để hiển thị nhiều thông tin hơn về cost và trace. Khi thêm các thành phần mới, cần cập nhật provenance, benchmark và các contract liên quan.

\subsection{Kết luận}
Qua đồ án, nhóm em xây dựng được một quy trình tìm đường gồm UCS, A*, các thuật toán tìm kiếm khác và hai phương pháp tối ưu tour. UCS được dùng làm mốc để kiểm tra A*. Trên 15 trường hợp kiểm thử lower-bound, A* cho cùng cost với UCS và không khám phá nhiều node hơn UCS. Kết quả này cho thấy heuristic đang dùng phù hợp với composite cost và các scenario đã kiểm tra.
